\documentclass{article}

\usepackage[a4paper, left=1.5cm, right=1.5cm, top=2cm, bottom=2cm]{geometry}
\usepackage{../../../../components/components}

\usepackage{fancyhdr}
\usepackage{hyperref}
\usepackage{ccicons}

% Configuration des en-têtes et pieds de page
\pagestyle{fancy}
\fancyhf{} % reset tout

\fancyhead[L]{Préparation CAPES}
\fancyhead[C]{Arithmétique}
\fancyhead[R]{2026-2027}

\fancyfoot[L]{Ewen Rodrigues de Oliveira\\\ccbyncsa}
\fancyfoot[R]{\thepage}

\begin{document}
\docTitle{Chapitre 1 : Construction de $\mathbb{N}$, $\mathbb{Z}$ et $\mathbb{Q}$}

\section{Construction de $\mathbb{N}$}

\subsection{Axiomes de Peano et opération successeur}
\definition{
    Les axiomes de Peano sont un ensemble d'axiomes qui définissent les propriétés fondamentales des nombres naturels.\\
    Ils sont les suivants :
    \begin{itemize}
        \item Il existe un nombre naturel $0$, qui n'est le successeur d'aucun nombre naturel.
        \item Chaque nombre naturel $n$ a un successeur, noté $S(n)$.
        \item Si $S(n) = S(m)$, alors $n = m$ (injectivité du successeur).
        \item Si une propriété est vraie pour $0$ et si elle est vraie pour le successeur de tout nombre naturel pour lequel elle est vraie, alors elle est vraie pour tous les nombres naturels (principe de récurrence).
    \end{itemize}
}

\theorem{Théorème}{}{true}{
    Les axiomes de Peano définissent une structure unique à isomorphisme près, appelée les nombres naturels $\mathbb{N}$.
}

\definition{
    L'opération successeur $S$ est définie de manière à ce que $S(n) = n + 1$ pour tout $n \in \mathbb{N}$.
}

\attention{Ici, $+$ et $1$ ne sont pas encore définis : ce sont des conventions d'écriture pour faciliter la compréhension, mais ils seront formellement définis plus tard à partir de l'opération successeur.}

\example{
    $S(0) = 1$, $S(1) = S(S(0)) = 2$, $S(2) = S(S(S(0))) = 3$, etc.
}

\remark{Ici, on remarque quelque chose d'important : les axiomes de Peano ne définissent pas explicitement l'opération d'addition ou de multiplication ni les relations d'ordre, mais ils permettent de les construire à partir de l'opération successeur.}

\subsection{Opération d'addition et de multiplication}
\subsubsection{Addition}
\definition{
    L'addition est définie de manière récursive à partir de l'opération successeur :
    \begin{itemize}
        \item $n + 0 = n$ pour tout $n \in \mathbb{N}$.
        \item $n + S(m) = S(n + m)$ pour tout $n, m \in \mathbb{N}$.
    \end{itemize}
}

\example{
    $2 + 3 = S(S(0)) + S(S(S(0))) = S(S(S(0) + S(S(0)))) = S(S(S(S(0) + S(0)))) = S(S(S(S(S(0) + 0)))) = S(S(S(S(S(0))))) = 5$.
}

\remark{Une telle définition de l'addition peut sembler compliquée, mais elle est essentielle pour construire rigoureusement les propriétés de l'addition à partir des axiomes de Peano. Néanmoins, on peut garder notre manière habituelle d'écrire $2 + 3$ pour faciliter la compréhension, tout en sachant que cela repose sur une construction formelle à partir de l'opération successeur.}

\theorem{Proposition}{}{true}{
    L'addition ainsi définie est associative, commutative et possède un élément neutre $0$.
}

\noindent{\textbf{Preuve :}}
\begin{itemize}
    \item Associativité : On montre que $(n + m) + k = n + (m + k)$ pour tous $n, m, k \in \mathbb{N}$.\\
    Par récurrence sur $k$ :
    \begin{itemize}
        \item $k = 0$ : $(n + m) + 0 = n + m = n + (m + 0)$.
        \item $k = S(l)$ : $(n + m) + S(l) = S((n + m) + l) = S(n + (m + l)) = n + S(m + l) = n + (m + S(l))$.
    \end{itemize}
    \item Commutativité : On montre que $n + m = m + n$ pour tous $n, m \in \mathbb{N}$.\\
    Par récurrence sur $m$ :
    \begin{itemize}
        \item $m = 0$ : $n + 0 = n = 0 + n$.
        \item $m = S(l)$ : $n + S(l) = S(n + l) = S(l + n) = S(l) + n$.
    \end{itemize}
    \item Élément neutre : $n + 0 = n$ pour tout $n \in \mathbb{N}$.
\end{itemize}

\subsubsection{Multiplication}
\definition{
    La multiplication est définie de manière récursive à partir de l'addition :
    \begin{itemize}
        \item $n \cdot 0 = 0$ pour tout $n \in \mathbb{N}$.
        \item $n \cdot S(m) = n \cdot m + n$ pour tout $n, m \in \mathbb{N}$.
    \end{itemize}
}

\example{
    $2 \cdot 3 = S(S(0)) \cdot S(S(S(0))) = S(S(0)) \cdot S(S(0)) + S(S(0)) = S(S(0)) \cdot S(0) + S(S(0)) + S(S(0)) = S(S(0)) \cdot 0 + S(S(0)) + S(S(0)) + S(S(0)) = 6$.
}

\theorem{Proposition}{}{true}{
    La multiplication ainsi définie est associative, commutative et possède un élément neutre $1$.
}

\noindent{\textbf{Preuve :}}
\begin{itemize}
    \item Associativité : On montre que $(n \cdot m) \cdot k = n \cdot (m \cdot k)$ pour tous $n, m, k \in \mathbb{N}$.\\
    Par récurrence sur $k$ :
    \begin{itemize}
        \item $k = 0$ : $(n \cdot m) \cdot 0 = 0 = n \cdot (m \cdot 0)$.   
        \item $k = S(l)$ : $(n \cdot m) \cdot S(l) = (n \cdot m) \cdot l + n \cdot m = n \cdot (m \cdot l) + n \cdot m = n \cdot (m \cdot S(l))$.
    \end{itemize}
    \item Commutativité : On montre que $n \cdot m = m \cdot n$ pour tous $n, m \in \mathbb{N}$.\\
    Par récurrence sur $m$ :
    \begin{itemize}
        \item $m = 0$ : $n \cdot 0 = 0 = 0 \cdot n$.
        \item $m = S(l)$ : $n \cdot S(l) = n \cdot l + n = l \cdot n + n = S(l) \cdot n$.
    \end{itemize}
    \item Élément neutre : $n \cdot 1 = n$ pour tout $n \in \mathbb{N}$.
\end{itemize}

\subsection{Relation d'ordre}

\definition{
    On définit la relation d'ordre $\leq$ sur $\mathbb{N}$ de la manière suivante : $n \leq m$ si et seulement s'il existe un $k \in \mathbb{N}$ tel que $n + k = m$.
}

\example{
    $2 \leq 5$ car $2 + 3 = 5$, mais $5 \nleq 2$ car il n'existe pas de $k \in \mathbb{N}$ tel que $5 + k = 2$.
}

\theorem{Proposition}{}{true}{
    La relation $\leq$ ainsi définie est un ordre total sur $\mathbb{N}$.
}
\noindent{\textbf{Preuve :}}\\
Pour qu'une relation soit un ordre total, elle doit être réflexive, antisymétrique, transitive et totale. On a :
\begin{itemize}
    \item Réflexivité : $n \leq n$ car $n + 0 = n$.
    \item Antisymétrie : Si $n \leq m$ et $m \leq n$, alors il existe $k, l \in \mathbb{N}$ tels que $n + k = m$ et $m + l = n$. En substituant $m$ dans la deuxième équation, on obtient $n + k + l = n$, ce qui implique que $k + l = 0$. Comme $k, l \in \mathbb{N}$, cela signifie que $k = 0$ et $l = 0$, donc $n = m$.
    \item Transitivité : Si $n \leq m$ et $m \leq p$, alors il existe $k, l \in \mathbb{N}$ tels que $n + k = m$ et $m + l = p$. En substituant $m$ dans la deuxième équation, on obtient $n + k + l = p$, ce qui implique que $n \leq p$.
    \item Totalité : Pour tous $n, m \in \mathbb{N}$, soit $n \leq m$ soit $m \leq n$. En effet, si $n \leq m$, alors il existe $k \in \mathbb{N}$ tel que $n + k = m$. Si $m \leq n$, alors il existe $l \in \mathbb{N}$ tel que $m + l = n$.
\end{itemize}
\section{Construction de $\mathbb{Z}$}

\subsection{Relation d'équivalence sur $\mathbb{N}^2$}

L'objectif de cette section est de construire un ensemble où toute soustraction devient possible. Pour ce faire, on formalise l'idée qu'un entier relatif peut être représenté par la différence de deux entiers naturels.

\definition{
    On définit sur l'ensemble des couples d'entiers naturels $\mathbb{N} \times \mathbb{N}$ la relation $\sim$ par :
    $$ (a, b) \sim (c, d) \iff a + d = b + c $$
    pour tous $a, b, c, d \in \mathbb{N}$.
}

\theorem{Proposition}{}{false}{
    La relation $\sim$ est une relation d'équivalence sur $\mathbb{N} \times \mathbb{N}$.
}

\noindent{\textbf{Preuve :}}
\begin{itemize}
    \item \textbf{Réflexivité :} Pour tout $(a, b) \in \mathbb{N}^2$, on a $a + b = b + a$ par commutativité de l'addition dans $\mathbb{N}$, donc $(a, b) \sim (a, b)$.
    \item \textbf{Symétrie :} Si $(a, b) \sim (c, d)$, alors $a + d = b + c$, ce qui implique $c + b = d + a$, d'où $(c, d) \sim (a, b)$.
    \item \textbf{Transitivité :} Si $(a, b) \sim (c, d)$ et $(c, d) \sim (e, f)$, alors $a + d = b + c$ et $c + f = d + e$.\\
    En sommant membre à membre, on obtient : $(a + d) + (c + f) = (b + c) + (d + e)$.\\
    En utilisant la commutativité et l'associativité dans $\mathbb{N}$, on peut simplifier par $c + d$ des deux côtés (propriété de régularité de l'addition dans $\mathbb{N}$), ce qui donne $a + f = b + e$. Par conséquent, $(a, b) \sim (e, f)$.
\end{itemize}

\definition{
    L'ensemble quotient $(\mathbb{N} \times \mathbb{N}) / \sim$ est noté $\mathbb{Z}$. C'est l'ensemble des entiers relatifs.\\
    La classe d'équivalence du couple $(a, b)$ est notée $\overline{(a, b)}$.
}

\remark{Intuitivement, la classe $\overline{(a, b)}$ représente l'entier relatif correspondant à la différence $a - b$.}

\example{
    \begin{itemize}
        \item La classe $\overline{(3, 1)} = \{(3, 1), (4, 2), (5, 3), \dots\}$ correspond à l'entier $2$.
        \item La classe $\overline{(0, 2)} = \{(0, 2), (1, 3), (2, 4), \dots\}$ correspond à l'entier $-2$.
        \item La classe $\overline{(1, 1)} = \{(0, 0), (1, 1), (2, 2), \dots\}$ correspond à l'entier $0$.
    \end{itemize}
}

\subsection{Opérations et structure d'anneau}

On munit désormais $\mathbb{Z}$ d'une loi d'addition et d'une loi de multiplication en prolongeant celles de $\mathbb{N}$.

\definition{
    Pour toutes classes $\overline{(a, b)}$ et $\overline{(c, d)}$ dans $\mathbb{Z}$, on définit :
    \begin{itemize}
        \item \textbf{L'addition :} $\overline{(a, b)} +_\mathbb{Z} \, \overline{(c, d)} = \overline{(a + c, b + d)}$
        \item \textbf{La multiplication :} $\overline{(a, b)} \cdot_{\mathbb{Z}} \, \overline{(c, d)} = \overline{(ac + bd, ad + bc)}$
    \end{itemize}
}

\attention{Il est fondamental de vérifier que ces opérations sont bien définies, c'est-à-dire que le résultat ne dépend pas des représentants choisis dans chaque classe d'équivalence.}

\theorem{Théorème}{}{true}{
    Muni de ces deux opérations, $(\mathbb{Z}, +, \cdot)$ est un anneau commutatif unitaire.
    \begin{itemize}
        \item L'élément neutre pour l'addition est $\overline{(0, 0)}$.
        \item L'opposé de la classe $\overline{(a, b)}$ est $-\overline{(a, b)} = \overline{(b, a)}$.
        \item L'élément neutre pour la multiplication est $\overline{(1, 0)}$.
    \end{itemize}
}

\subsection{Plongement de $\mathbb{N}$ dans $\mathbb{Z}$ et notation usuelle}

\theorem{Proposition}{}{true}{
    L'application $\phi : \mathbb{N} \to \mathbb{Z}$ définie par $\phi(n) = \overline{(n, 0)}$ est un morphisme injectif d'anneaux non unitaires (il préserve l'addition, la multiplication, et l'ordre).
}

\remark{Grâce à cette injection, on identifie chaque entier naturel $n$ à sa classe $\overline{(n, 0)}$. On peut ainsi écrire $\mathbb{N} \subset \mathbb{Z}$.}

\vocabulary{
    Par convention et pour simplifier l'écriture :
    \begin{itemize}
        \item Si $a \geq b$, il existe $k \in \mathbb{N}$ tel que $a = b + k$. Alors $\overline{(a, b)} = \overline{(b+k, b)} = \overline{(k, 0)}$, que l'on note simplement $k$.
        \item Si $a < b$, il existe $k \in \mathbb{N}^*$ tel que $b = a + k$. Alors $\overline{(a, b)} = \overline{(a, a+k)} = \overline{(0, k)}$, que l'on note simplement $-k$.
    \end{itemize}
}

\subsection{Relation d'ordre sur $\mathbb{Z}$}

\definition{
    On définit la relation d'ordre $\leq$ sur $\mathbb{Z}$ par :
    $$ \overline{(a, b)} \leq \overline{(c, d)} \iff a + d \leq b + c \quad (\text{dans } \mathbb{N}) $$
}

\theorem{Proposition}{}{true}{
    La relation $\leq$ est un ordre total sur $\mathbb{Z}$, compatible avec l'addition et la multiplication par des éléments positifs.
}

\section{Construction de $\mathbb{Q}$}

\subsection{Relation d'équivalence sur $\mathbb{Z} \times \mathbb{Z}^*$}

L'objectif de cette section est de construire un corps où toute équation de la forme $mx = n$ (avec $m \neq 0$) possède une solution unique. On formalise ici l'idée qu'un nombre rationnel peut être représenté par le quotient d'un entier relatif par un entier relatif non nul.

\definition{
    On note $\mathbb{Z}^* = \mathbb{Z} \setminus \{\overline{(0,0)}\}$ l'ensemble des entiers relatifs non nuls. On définit sur l'ensemble des couples $\mathbb{Z} \times \mathbb{Z}^*$ la relation $\mathcal{R}$ par :
    $$ (a, b) \mathcal{R} (c, d) \iff a \cdot_{\mathbb{Z}} d = b \cdot_{\mathbb{Z}} c $$
    pour tous $a, c \in \mathbb{Z}$ et $b, d \in \mathbb{Z}^*$.
}

\theorem{Proposition}{}{false}{
    La relation $\mathcal{R}$ est une relation d'équivalence sur $\mathbb{Z} \times \mathbb{Z}^*$.
}

\noindent{\textbf{Preuve :}}
\begin{itemize}
    \item \textbf{Réflexivité :} Pour tout $(a, b) \in \mathbb{Z} \times \mathbb{Z}^*$, on a $a \cdot_{\mathbb{Z}} b = b \cdot_{\mathbb{Z}} a$ par commutativité de la multiplication dans $\mathbb{Z}$, donc $(a, b) \mathcal{R} (a, b)$.
    \item \textbf{Symétrie :} Si $(a, b) \mathcal{R} (c, d)$, alors $a \cdot_{\mathbb{Z}} d = b \cdot_{\mathbb{Z}} c$, ce qui implique $c \cdot_{\mathbb{Z}} b = d \cdot_{\mathbb{Z}} a$ par commutativité, d'où $(c, d) \mathcal{R} (a, b)$.
    \item \textbf{Transitivité :} Si $(a, b) \mathcal{R} (c, d)$ et $(c, d) \mathcal{R} (e, f)$, alors $a \cdot_{\mathbb{Z}} d = b \cdot_{\mathbb{Z}} c$ et $c \cdot_{\mathbb{Z}} f = d \cdot_{\mathbb{Z}} e$.\\
    Multiplions la première égalité par $f$ et la seconde par $b$ :
    $$ (a \cdot_{\mathbb{Z}} d) \cdot_{\mathbb{Z}} f = (b \cdot_{\mathbb{Z}} c) \cdot_{\mathbb{Z}} f \quad \text{et} \quad (c \cdot_{\mathbb{Z}} f) \cdot_{\mathbb{Z}} b = (d \cdot_{\mathbb{Z}} e) \cdot_{\mathbb{Z}} b $$
    En combinant ces deux égalités par transitivité dans $\mathbb{Z}$, on obtient :
    $$ (a \cdot_{\mathbb{Z}} f) \cdot_{\mathbb{Z}} d = (e \cdot_{\mathbb{Z}} b) \cdot_{\mathbb{Z}} d $$
    Puisque $\mathbb{Z}$ est un anneau intègre et que $d \in \mathbb{Z}^*$ ($d$ est non nul), on peut simplifier par $d$ (propriété de régularité), ce qui donne $a \cdot_{\mathbb{Z}} f = b \cdot_{\mathbb{Z}} e$. Par conséquent, $(a, b) \mathcal{R} (e, f)$.
\end{itemize}

\definition{
    L'ensemble quotient $(\mathbb{Z} \times \mathbb{Z}^*) / \mathcal{R}$ est noté $\mathbb{Q}$. C'est l'ensemble des nombres rationnels.\\
    La classe d'équivalence du couple $(a, b)$ est notée $\frac{a}{b}$ ou $a/b$.
}

\remark{Intuitivement, la classe $\frac{a}{b}$ formalise le quotient algébrique de l'entier $a$ par l'entier non nul $b$.}

\example{
    \begin{itemize}
        \item La classe $\frac{1}{2} = \{(1, 2), (2, 4), (-3, -6), \dots\}$ correspond au rationnel un demi.
        \item La classe $\frac{-2}{3} = \{(-2, 3), (2, -3), (-4, 6), \dots\}$ correspond au rationnel moins deux tiers.
        \item La classe $\frac{0}{1} = \{(0, 1), (0, 2), (0, -5), \dots\}$ correspond au rationnel nul.
    \end{itemize}
}

\subsection{Opérations et structure de corps}

On définit sur $\mathbb{Q}$ une loi d'addition et une loi de multiplication naturelles induites par les opérations de l'anneau $\mathbb{Z}$.

\definition{
    Pour toutes classes $\frac{a}{b}$ et $\frac{c}{d}$ dans $\mathbb{Q}$, on définit :
    \begin{itemize}
        \item \textbf{L'addition :} $\frac{a}{b} +_\mathbb{Q} \frac{c}{d} = \frac{a \cdot_{\mathbb{Z}} d +_{\mathbb{Z}} b \cdot_{\mathbb{Z}} c}{b \cdot_{\mathbb{Z}} d}$
        \item \textbf{La multiplication :} $\frac{a}{b} \cdot_{\mathbb{Q}} \frac{c}{d} = \frac{a \cdot_{\mathbb{Z}} c}{b \cdot_{\mathbb{Z}} d}$
    \end{itemize}
}

\attention{Comme $b \neq 0$ et $d \neq 0$, par intégrité de l'anneau $\mathbb{Z}$, le produit $b \cdot_{\mathbb{Z}} d \neq 0$, le représentant appartient donc bien à $\mathbb{Z} \times \mathbb{Z}^*$. De plus, il convient de vérifier que ces lois sont indépendantes du choix des représentants choisis dans chaque classe.}

\theorem{Théorème}{}{true}{
    Muni de ces deux opérations, $(\mathbb{Q}, +, \cdot)$ est un corps commutatif.
    \begin{itemize}
        \item L'élément neutre pour l'addition est la classe $\frac{0}{1}$.
        \item L'opposé de la classe $\frac{a}{b}$ est $-\frac{a}{b} = \frac{-a}{b} = \frac{a}{-b}$.
        \item L'élément neutre pour la multiplication est la classe $\frac{1}{1}$.
        \item Tout élément non nul $\frac{a}{b} \in \mathbb{Q}^*$ (ce qui implique $a \neq 0$) admet un inverse unique donné par $\left(\frac{a}{b}\right)^{-1} = \frac{b}{a}$.
    \end{itemize}
}

\subsection{Plongement de $\mathbb{Z}$ dans $\mathbb{Q}$ et notations usuelles}

\theorem{Proposition}{}{true}{
    L'application $\psi : \mathbb{Z} \to \mathbb{Q}$ définie par $\psi(a) = \frac{a}{1}$ est un morphisme injectif de corps (il préserve l'identité de l'addition, de la multiplication et de la relation d'ordre).
}

\remark{Grâce à cette injection canonique, on identifie chaque entier relatif $a$ à la classe $\frac{a}{1}$. On pose ainsi légitimement l'inclusion d'anneaux $\mathbb{Z} \subset \mathbb{Q}$.}

\vocabulary{
    Tout nombre rationnel peut s'écrire sous la forme unique d'une fraction irréductible $\frac{p}{q}$ où $p \in \mathbb{Z}$, $q \in \mathbb{N}^*$ et $\text{pgcd}(p, q) = 1$.
}

\subsection{Relation d'ordre sur $\mathbb{Q}$}

Pour préserver la compatibilité de l'ordre avec la multiplication par un élément positif, on s'assure d'abord de travailler avec des dénominateurs strictement positifs.

\definition{
    Soient deux rationnels exprimés avec un dénominateur positif (quitte à multiplier numérateur et dénominateur par $-1$), c'est-à-dire $\frac{a}{b}$ et $\frac{c}{d}$ avec $b > 0$ et $d > 0$ dans $\mathbb{Z}$. On définit la relation d'ordre $\leq$ sur $\mathbb{Q}$ par :
    $$ \frac{a}{b} \leq \frac{c}{d} \iff a \cdot_{\mathbb{Z}} d \leq b \cdot_{\mathbb{Z}} c \quad (\text{dans } \mathbb{Z}) $$
}

\theorem{Proposition}{}{true}{
    La relation $\leq$ est un ordre total sur $\mathbb{Q}$, faisant de $(\mathbb{Q}, +, \cdot, \leq)$ un corps ordonné archimédien.
}

\newpage
\tableofcontents

\section*{Crédits}
Ce cours est mis à disposition selon les termes de la Licence Creative Commons \ccbyncsa\ \textbf{Attribution - Pas d'Utilisation Commerciale - Partage dans les Mêmes Conditions 4.0 International}. 
Pour voir une copie de cette licence, visitez \url{http://creativecommons.org/licenses/by-nc-sa/4.0/}.

\end{document}
